\section{Investigación de necesidades y antecedentes}

\subsection{Digitalización de pequeñas y medianas empresas}

El Plan Nacional de Desarrollo e Inversión Pública incluye acciones para
promover la digitalización de pymes mediante plataformas de comercialización y
herramientas de trabajo \parencite{meic_pndip2023_2026}. En Pulso, este contexto
orienta la decisión de reunir primero pagos, membresías, asistencia, planes y
progreso. Así, el gimnasio puede aprender a usar la plataforma por etapas.

\subsection{Situación del mercado y la industria}

La Encuesta Nacional de Microempresas de los Hogares 2024 estimó 435\,779
microempresas en Costa Rica y registró un crecimiento de 10,5\% frente a 2023
\parencite{inec_enameh2024}. Ese universo nacional no equivale al número de
gimnasios ni permite calcular la demanda de Pulso, pero confirma que los
negocios pequeños constituyen una parte relevante del entorno productivo. El
Plan Nacional de Desarrollo e Inversión Pública 2023--2026 incluye entre sus
acciones promover la digitalización de las pymes mediante plataformas de
comercialización y herramientas de trabajo \parencite{meic_pndip2023_2026}.

La exploración de directorios públicos permitió ubicar posibles establecimientos
y centros de entrenamiento asociados con Pérez Zeledón, aunque sus listados no
constituyen un censo de gimnasios activos y pueden incluir registros
desactualizados o servicios distintos \parencite{cybo_gimnasios_san_isidro}. Esta
base sirve para iniciar la prospección local, no para afirmar un tamaño de
mercado. Además, cuatro de los cinco responsables consultados señalaron
dificultades relacionadas con los cobros y los pagos pendientes, y cuatro
manifestaron interés en conocer una demostración antes de definir un precio
\parencite{encuesta_gymhub_2026}. Estos resultados no garantizan contrataciones,
pero sí muestran una necesidad reconocida y una oportunidad comercial que
merece validarse mediante un piloto. Por ello, Pulso puede iniciar con
establecimientos identificados, demostrar cómo la plataforma facilita la
gestión del gimnasio y medir de forma directa cuántos avanzan a una contratación.
El crecimiento proyectado se considerará una meta por validar, no una venta
asegurada.

\subsection{Origen de la propuesta}

El equipo está integrado por dos estudiantes de undécimo año de la especialidad
Configuración y Soporte a Redes de Computadoras y Sistemas Operativos de Red. La
idea, denominada inicialmente GymHub y hoy identificada comercialmente como
Pulso, surgió en
décimo año, durante el segundo periodo de 2025, como proyecto de la subárea
Fundamentos de Programación impartida por el profesor Andrés Mena Abarca. Al
finalizar ese periodo ya se contaba con la propuesta, los requisitos principales
y una primera versión del sistema.

Durante 2026, aquella primera versión se amplió hasta convertirla en un
sistema integral en línea con módulos de miembros, planes de entrenamiento, asistencia,
progreso y facturación
\parencite{gymhub_repo}. Esta es la primera participación de Pulso en
ExpoTÉCNICA; por lo tanto, es un proyecto nuevo para la feria y no un proyecto de
continuación.

La encuesta complementó el desarrollo técnico con información de posibles
clientes y personas usuarias. Los diez miembros calificaron con 4 o 5 la utilidad
de reunir pagos, asistencia, rutina y progreso. Además, cuatro de cinco
responsables identificaron los cobros y pagos pendientes como una dificultad.
Estas respuestas llevaron al equipo a dar prioridad a los cobros, las rutinas
y la consulta de información desde el celular durante la demostración.

\subsection{Alternativas actuales y competencia}

Antes de comparar Pulso con proveedores específicos, se consideraron las
alternativas que ya utilizan los gimnasios sin adquirir una plataforma integral.
La Tabla~\ref{tab:alternativas} resume sus ventajas y la diferencia que ofrece la
propuesta.

\begin{table}[htbp]
\caption{Comparación preliminar con alternativas de gestión}
\label{tab:alternativas}
\begin{tabularx}{\textwidth}{>{\bfseries}p{2.8cm}X X}
\toprule
Alternativa & Ventaja & Limitación frente a la propuesta \\
\midrule
Cuaderno o archivo físico & Bajo costo y aprendizaje inmediato. & Riesgo de deterioro, búsqueda manual y dificultad para relacionar registros. \\
Hoja de cálculo & Flexible y conocida. & Depende del diseño del encargado y no ofrece por sí sola experiencia separada por rol. \\
Mensajería & Comunicación directa y uso extendido. & Las conversaciones no funcionan como base estructurada de pagos, membresías o asistencia. \\
Aplicaciones separadas & Permiten resolver una tarea concreta. & Pueden fragmentar datos y aumentar cambios de herramienta. \\
Pulso & Integra operación administrativa y seguimiento deportivo. & Requiere conexión, capacitación, soporte y validación comercial. \\
\bottomrule
\end{tabularx}
\caption*{\textit{Nota}. Elaboración propia. La encuesta confirmó el uso combinado de estas alternativas en cinco casos administrativos, 2026.}
\end{table}

La comparación parte de las herramientas mencionadas en la encuesta. Pulso
busca facilitar la consulta conjunta de pagos, membresías y entrenamiento, con
ayuda para empezar a usarlo. Su alcance inicial no incluye procesamiento
bancario, contabilidad, reservas ni apertura automática de puertas.

\subsection{Criterios de diseño derivados}

El equipo definió cinco criterios de diseño: facilidad de uso, acceso desde
celular, roles claros, aprendizaje gradual y costos que se puedan cubrir.
También se tomó en cuenta el manejo de información personal. La Ley n.º 8968
protege la autodeterminación informativa y se aplica al tratamiento automatizado
o manual de datos personales en bases públicas o privadas
\parencite{ley8968}. Por ello, privacidad y consentimiento forman parte de la
preparación del servicio desde el inicio. Los cinco
responsables calificaron el control de acceso como importante o muy importante.
Además, nueve de diez miembros indicaron el celular como dispositivo principal,
lo que refuerza la prioridad de pantallas adaptables al celular. La pregunta de
privacidad no se aplicó a los miembros; ese tema solo se analizó entre los
responsables.
